\section{DISCUSSION}
\subsection{Shape of drying curves compared to theory}
\begin{itemize}
    \item \textbf{Warming-up}: Not clearly visible in experimental curves because it occurs very fast (ignored recordings usually start at 2-3 mins).
    \item \textbf{Constant vs. Falling}: Visible, though transition points (critical moisture $U_{th}$) are sensitive to sampling frequency.
    \item \textbf{Temperature Effect}: Consistent with theory; higher temperatures yield higher drying rates (lower curves for lower temperatures).
\end{itemize}

\subsection{Analysis of results}
\begin{itemize}
    \item \textbf{Constant Rate $N$}: Matches theory conceptually; higher $T$ yields higher $N$.
    \item \textbf{Relative drying coefficient $\chi$}: Theoretically only dependent on material, but experimental values vary because they rely on calculated $U_{th}$ and $U^*$.
    \item \textbf{Drying coefficient $K$}: Suitable trend; $K$ increases with temperature as $N$ increases.
    \item \textbf{Drying time $\tau$}: $\tau_1$ and $\tau_2$ generally decrease with increasing temperature, though experimental errors in $N$ affect $\tau_1$ suitability.
\end{itemize}

\subsection{Error evaluation}
\begin{itemize}
    \item \textbf{Critical moisture $U_{th}$}: Large errors ($39\% \div 87\%$) due to plotting approximations.
    \item \textbf{Constant rate $N$}: Very large errors ($130\% \div 217\%$). Causes include:
    \begin{itemize}
        \item Assuming constant $\alpha_p$ (empirical formula $\alpha_p = 0.0229 + 0.0174v_k$).
        \item Inaccuracies in determining $P_m$ and $P$ from graphs.
        \item Manual plotting of the drying curve differentials.
    \end{itemize}
    \item \textbf{Time $\tau$}: Negative errors in $\tau_1$ (Exp < Theory) because experimental $N$ was significantly larger than theoretical predictions.
\end{itemize}

\subsection{Methods for remediation}
\begin{itemize}
    \item \textbf{Experiment}: Ensure materials are 100\% dry for $G_c$, use electronic sensors for precise recording.
    \item \textbf{Calculation}: Use formulas instead of manual graph lookups, apply minimum mean square error for plotting.
\end{itemize}
